The consolidation--disruption (CD) index is used in research evaluation, but how it scores papers later retracted is
unclear. On OpenCitations linked to PubMed, we compared late-retracted biomedical papers with never-retracted controls
from the same journal and year. Late-retracted means that the first retraction notice occurred at least six calendar
years after publication. The pre-specified five-year index (CD$_5$) showed a small mean gap (0.0002) and no clear
top-tier excess, but positive scores were more common (odds ratio 1.23, 95\% confidence interval 1.13--1.34).
An exploratory analysis, planned after these results were known, found that papers citing a retracted paper were
less likely to also cite its references: on average, 65.5\% versus 68.8\%.
CD$_5$ compresses this contrast because papers citing only the references dominate its denominator.
The variant without that term (DI$_{nok}$) had a mean gap of 0.067 (0.034--0.095).
As a robustness check, excluding notice, comment/letter and retracted citing records increased the gap to 0.082.
The filtered contrast persisted with reference and citation adjustment
and after excluding the ten journals with the most late-retracted papers in the study.
For unfiltered scores, a small topic-matched subset did not establish a gap.
CD$_5$ can look neutral while DI$_{nok}$ scores late-retracted papers higher within journal/year comparisons.
Researchers and bulk score providers should end scoring windows before the first retraction notice,
with notice records removed.
